\section{Appendix and Notes}\label{sec:appendix}

\subsection{Analysis script}

The identification was carried out in MATLAB by \texttt{simulate\_lab1.m}, submitted with this
report and run as \texttt{matlab -batch simulate\_lab1}. It takes no arguments and resolves
every path against its own location, so it runs from any working directory. It reads the logs
from \texttt{data/lab-session-2026-08-31/} beside it and writes the figures of this report to
\texttt{report/figures/}, creating that folder if it is not there. The sine files are named
\texttt{sine\_<frequency>.csv} so the drive frequency is read off the file name. A trailing \texttt{\_a3} marks a run driven at amplitude
$3$, held out of the frequency points and used for the linearity check instead.
Listing~\ref{lst:matlab_out} is the output it produced. Every measured quantity in this report
is printed there, including the readings taken off the Bode plot by hand. Those are the recovered DC
gain, both versions of the resonant peak, the phase crossings, the roll-off slopes and the
excess phase per unit frequency.

What is not printed is the arithmetic done on those numbers in the text, such as the error
bands carried through Equation~\ref{eqn:logdec}, the counterfactual readings used to show
what a misread peak would have cost, and the resolution figures quoted in millimetres. Those
are single steps from values that are printed and are shown in place where they are used.

It reads the gain from the change in the final value, the damping ratio from the first
overshoot with the logarithmic decrement as a cross-check, and the damped period from the
peaks in the record. Only peaks that clear the final value by at least three times the
output resolution are counted, since a sample sitting one resolution step above the settled
value is quantisation noise, not a genuine overshoot, and counting it drags the
averaged period badly. Where only one peak survives that test, the damped period is taken as
twice the time to that peak. The logarithmic decrement is read off the first peak and the
first trough rather than off two peaks, since successive extrema alternate in sign and the
trough is one decay factor away from the peak while the second peak is two. That is what
makes the cross-check available at all on this data, where the second peak has already
fallen under the resolution. Amplitude and phase at each drive frequency are fitted by least
squares against a sinusoid at the known frequency, which avoids reading peaks off records
sampled at 100~ms.

The frequency fits are Gauss-Newton with a numerical Jacobian, started from the step answer,
and the script runs three. One frees all four parameters and is the independent reading of the
plant from the frequency points quoted in Section~\ref{sec:system_id}. One repeats that on the
seven points at or below $2$~rad/s, which is where the gain and the input scaling are
confirmed. The third holds the delay fixed and fits the remaining three around it, at each of
the three delays of Table~\ref{tab:refit}. The real against complex pole check is also a search,
but over a grid of rate pairs with the linear part solved exactly at each. The script also refits every step log in turn, which
is where the repeatability spread comes from, and compares the $1$~rad/s sinusoid against its
amplitude $3$ repeat for the linearity check.

The script stops with an error if no peak in the step record clears the resolution, since the
formulae used here are the underdamped ones and an overdamped record would need a different
reading. That is not exercised on this data.

\subsection{MATLAB output}

Listing~\ref{lst:matlab_out} is what \texttt{simulate\_lab1.m} prints when it is run on the
logs, reproduced as it appears in the console. Every identified value quoted in
Section~\ref{sec:system_id} can be read off it, and the measured against model columns of the
frequency table are the numbers plotted in Figure~\ref{fig:validation_freq}.

\lstinputlisting[numbers=none,
                 caption={Console output of \texttt{simulate\_lab1.m} run on the logs of
                 31 August 2026.},
                 label=lst:matlab_out]{matlab_output.txt}

\texttt{simulate\_lab1.m} is submitted alongside this report rather than printed in it, since
it runs to several hundred lines and the console output above is what the numbers are read
from. Given the logs it repeats the whole identification. It reads the CSV files, identifies
$A$, $\zeta$ and $\omega_n$ from the step record, fits the transport delay to the phase left
over from the second order model, builds the transfer function of Equation~\ref{eqn:identified}
and reports how far the simulation sits from each experiment, $0.24$ dB RMS on magnitude and
$25.5$ degrees RMS on phase falling to $0.2$ degrees once the delay is included.

\subsection{Raw data}

All logs were captured on 31 August 2026 at a 100~ms sample period, and each run was copied
out and renamed before the next was started, since the simulator appends to a single fixed
file name.

\begin{itemize}
    \item \texttt{step1.csv}, \texttt{step2.csv}, \texttt{step3.csv}, a step of $1$~V applied
          through \textbf{Offset Input}, with all sinusoid fields zero. Three repeats of the
          same test, logged for $30$, $28$ and $41$~s after the step.
    \item \texttt{step\_amp2.csv}, the same test with a step of $2$~V, used to check that the
          response scales linearly with the input.
    \item \texttt{sine\_0p15.csv} through \texttt{sine\_10p0.csv}, twelve single frequency
          runs at $0.15$, $0.30$, $0.50$, $0.80$, $1.00$, $1.30$, $2.00$, $3.00$, $4.00$,
          $6.00$, $8.00$ and $10.0$~rad/s, each at drive amplitude $1$ with zero offset. The
          three runs below $0.8$~rad/s hold just over three complete cycles and the rest at
          least five, which leaves the fitted second half of every record well clear of the
          transient. The file name carries the drive frequency, with \texttt{p} standing in
          for the decimal point.
    \item \texttt{sine\_1p0\_a3.csv}, the $1.0$~rad/s run repeated at drive amplitude $3$,
          used as the frequency domain counterpart of the linearity check.
\end{itemize}
